\section{Contexte initial et objectifs du projet}
\subsection{Limites de l'outil existant}
% Description de la situation de départ
% Problèmes et inefficacités identifiés
% Impact sur les utilisateurs et l'entreprise

\subsection{Vision et objectifs de la refonte}
% Objectifs stratégiques et opérationnels
% Gains attendus
% Indicateurs de succès définis

\section{Périmètre et organisation du projet}
\subsection{Couverture fonctionnelle et technique}
% Périmètre géographique (25+ DR)
% Périmètre fonctionnel
% Profils utilisateurs concernés

\subsection{Gouvernance et méthodologie}
% Organisation de l'équipe projet
% Méthodologie appliquée (Agile, GitLab)
% Rôles et responsabilités

\section{Étapes clés et réalisations}
\subsection{Phase de cadrage et d'analyse}
% Spécification du besoin
% Ateliers avec les utilisateurs
% Formalisation des exigences

\subsection{Conception et développement}
% Architecture et choix techniques
% Sprints de développement
% Tests et validation

\subsection{Déploiement et adoption}
% Stratégie de déploiement
% Formation et accompagnement
% Gestion du changement

\section{Solutions techniques implémentées}
\subsection{Frontend et expérience utilisateur}
% Interface Angular
% Fonctionnalités principales
% Optimisations UI/UX

\subsection{Backend et intégration}
% Architecture NestJS
% Intégration avec le SI Enedis (SSO)
% Gestion des données avec Prisma

\subsection{Automatisation et traitement des données}
% Flux de données automatisés
% Imports mensuels
% Validation et intégrité des données

\section{Résultats et impacts}
\subsection{Métriques et bénéfices quantifiables}
% Réduction du temps de traitement (50%)
% Gain de productivité (1 jour/mois/DR)
% Autres indicateurs pertinents

\subsection{Bénéfices qualitatifs}
% Amélioration de l'expérience utilisateur
% Harmonisation des processus
% Fiabilité et qualité des données